\chapter{Stage-Gate and Development Strategy}

\section{Session Overview}

This session asks your group to decide how the project should be managed before detailed design work begins. In practice, a product team does not only choose \emph{what} to develop; it also decides \emph{how} the work will be organised, reviewed, and coordinated. These choices affect risk, speed, communication, and the quality of later design decisions.

By the end of this session, your group should be able to:
\begin{enumerate}
    \item identify a suitable product development process for your hydroponics project;
    \item select a development process flow that matches the project's complexity and uncertainty;
    \item distinguish between departmental structure and organisational structure; and
    \item justify a recommended team structure for the project.
\end{enumerate}

\section{Why a Stage-Gate Approach Matters}

A \textbf{stage-gate process} breaks product development into major stages separated by review points, or \emph{gates}. At each gate, the team checks whether the project should continue, be revised, or stop. This approach is useful because it forces teams to review evidence before moving forward.

For APID, the stage-gate mindset helps your group do three things:
\begin{itemize}
    \item confirm that the chosen opportunity still fits user needs and project goals;
    \item identify technical or market risks early; and
    \item coordinate decisions across engineering, design, communication, and planning tasks.
\end{itemize}

\section{Key Terms}

This section separates terms that are often confused.

\begin{itemize}
    \item \textbf{Product development process}: the overall type of development strategy, such as market-pull or technology-push.
    \item \textbf{Development process flow}: the sequence pattern used to execute the work, such as a generic linear flow or a spiral cycle.
    \item \textbf{Departmental structure}: the functions that must contribute to the work, such as marketing, engineering, manufacturing, or quality.
    \item \textbf{Organisational structure}: how authority and responsibility are arranged, such as a functional team, project team, lightweight structure, or heavyweight structure.
\end{itemize}

\section{Product Development Process}

Table~\ref{tab:session2-process-types} summarises common process types of Product Development Processes with Their
Features and Examples. When your group chooses a process type, focus on the product's dominant challenge. For example, if the main uncertainty is a novel sensor technology, a technology-push or high-risk framing may fit better than a generic market-pull process.

\begin{longtable}{|p{3.2cm}|p{4.2cm}|p{5.4cm}|p{2.5cm}|}
\caption{Types of Product Development Processes with Their Features and Examples}
\label{tab:session2-process-types} \\
\hline
\textbf{Process Type} & \textbf{Description} & \textbf{Distinct Features} & \textbf{Examples} \\
\hline
Generic (Market-Pull) Products & The team begins with a market opportunity and selects appropriate technologies to meet customer needs. & Process generally includes distinct planning, concept development, system-level design, detail design, testing and refinement, and production ramp-up phases. & Sporting goods, furniture, tools. \\
\hline
Technology-Push Products & The team begins with a new technology, then finds an appropriate market. & Planning phase involves matching technology and market. Concept development assumes a given technology. & Gore-Tex rainwear, Tyvek envelopes. \\
\hline
Platform Products & The team assumes that the new product will be built around an established technological subsystem. & Concept development assumes a proven technology platform. & Consumer electronics, computers, printers. \\
\hline
Process-Intensive Products & Characteristics of the product are highly constrained by the production process. & Either an existing production process must be specified from the start, or both product and process must be developed together from the start. & Snack foods, breakfast cereals, chemicals, semiconductors. \\
\hline
Customized Products & New products are slight variations of existing configurations. & Similarity of projects allows for a streamlined and highly structured development process. & Motors, switches, batteries, containers. \\
\hline
High-Risk Products & Technical or market uncertainties create high risks of failure. & Risks are identified early and tracked throughout the process. Analysis and testing activities take place as early as possible. & Pharmaceuticals, space systems. \\
\hline
Quick-Build Products & Rapid modeling and prototyping enables many design-build-test cycles. & Detail design and testing phases are repeated a number of times until the product is completed or time/budget runs out. & Software, cellular phones. \\
\hline
Product-Service Systems & Products and their associated service elements are developed simultaneously. & Both physical and operational elements are developed, with particular attention to design of the customer experience and the process flow. & Restaurants, software applications, financial services. \\
\hline
Complex Systems & System must be decomposed into several subsystems and many components. & Subsystems and components are developed by many teams working in parallel, followed by system integration and validation. & Airplanes, jet engines, automobiles. \\
\hline
\end{longtable}




\section{Product Development Process Flows}

After selecting the broad process type, choose the flow that best suits the project execution.

There are several common flow patterns:
\begin{itemize}
    \item \textbf{Generic flow}: a mostly sequential path from planning to concept development, design, testing, and refinement. This works well for low-to-moderate uncertainty.
    \item \textbf{Spiral flow}: repeated cycles of planning, prototyping, review, and revision. This is useful when technical learning must happen through iteration.
    \item \textbf{Complex-system flow}: parallel work across multiple subsystems followed by integration. This is appropriate when the product contains many interacting modules.
\end{itemize}

Analyze the characteristics and requirements of your team's selected product.Consider factors such as complexity, risk, and flexibility in the development process. Evaluate each Development Process Flow in relation to your product's needs and assess how each process flow aligns with the specific requirements and constraints of your team project.  Select the most suitable development process flow, for your team.


% If your product includes electronics, fluid handling, and structure that must work together, a spiral or complex-system flow is often more realistic than a single-pass linear flow.

\section{Departmental Structure}

This section focuses on \emph{who} needs to contribute. A product development team commonly draws from:
\begin{itemize}
    \item Marketing and sales;
    \item Engineering;
    \item Quality assurance
    \item Manufacturing;
    \item Purchasing;
    \item Legal
    \item Financial; and
    \item Project management.
\end{itemize}

You can refer to Exhibit 1-2 (Development of an electromechanical product) or Exhibit 2.7 (Tyco International) explore the different departmental structure during the rally point various types of product development projects.


For your own product development project, assess whether it is necessary to maintain, reduce, or add new departments to this list. Provide a well-reasoned explanation for your chosen approach.




\section{Organisational Structure}

The objective of this section in the lab manual is to guide you in understanding various organizational structures and selecting the most appropriate one for your specific product development project. By familiarizing yourself with these structures and making an informed decision, you will be able to effectively organize and manage your project team.

Begin by familiarizing yourself with different type of organizational structures which include 1) functional organization, 2) project organization, 3) lightweight organization, or 4) heavyweight project organization.

Determine the most appropriate organizational structure for your specific product development project. Take into account the choice you made in Activity 1 and Activity 2 from week 2.


\section{Pre-Class Preparation}

Before class, each group member should prepare short notes covering the following decisions:
\begin{enumerate}
    \item the most suitable development process type for the selected hydroponic product;
    \item the most suitable development process flow;
    \item the most important departments or functions needed for the project; and
    \item the organisational structure that would best coordinate the work.
\end{enumerate}

For each decision, write one short justification based on project complexity, uncertainty, user needs, technical risks, or team coordination.

\section{In-Class Activity: Consolidate the Development Strategy}

In class, compare the individual proposals and agree on one team recommendation for each category.

\subsection{Suggested Decision Steps}
\begin{enumerate}
    \item compare each member's proposed process type and identify the strongest justification;
    \item agree on the development process flow that best matches the level of iteration or subsystem integration required;
    \item decide which departments or functions are essential to the project;
    \item choose the organisational structure that supports those functions most effectively; and
    \item document the final decision in a short report for ITEL submission if requested.
\end{enumerate}

\subsection{Decision Matrix Template}

Use Table~\ref{tab:session2-decision-matrix} if your group wants a structured way to compare options.

\begin{table}[htbp]
\centering
\caption{Simple decision matrix for Session 2 selections}
\label{tab:session2-decision-matrix}
\renewcommand{\arraystretch}{1.2}
\begin{tabularx}{\textwidth}{|>{\raggedright\arraybackslash}p{0.22\textwidth}|>{\raggedright\arraybackslash}p{0.18\textwidth}|>{\raggedright\arraybackslash}p{0.16\textwidth}|X|}
\hline
\textbf{Decision area} & \textbf{Selected option} & \textbf{Confidence (1--5)} & \textbf{Justification} \\
\hline
Development process type &  &  &  \\
\hline
Development process flow &  &  &  \\
\hline
Departmental structure &  &  &  \\
\hline
Organisational structure &  &  &  \\
\hline
\end{tabularx}
\end{table}

\section{Summary and Next Steps}

By the end of this session, your group should have one agreed development strategy covering process type, process flow, departmental structure, and organisational structure. Keep this summary because it will influence later decisions about planning, customer needs, concept development, and project communication.
